% Chapter 4 --- Design Methodology and Iterative Prototyping

\chapter{Design Methodology and Iterative Prototyping}\doublespacing
\label{Chapter4}
\lhead{Chapter III. \emph{Design Methodology and Iterative Prototyping}}

%----------------------------------------------------------------------------------------
\section{Preliminary Visual Analysis and Colour Theory}
%----------------------------------------------------------------------------------------

A structured visual analysis was conducted across all twelve target categories using physical grain specimens from procurement sites. The working question was: \textit{which visual cues do experienced assessors use when making a category call, and can controlled illumination make those cues measurable by a camera?}

Early experiments used coloured chart-paper backgrounds at varied light angles. Warm amber tones improved contrast between \textit{Lokwan} and \textit{Kernel Burnt}; cooler backgrounds improved boundary definition for segmentation. No single background setting performed consistently across all twelve categories and wheat varieties.

This led to the first key design decision: replace static backgrounds with a programmable active substrate. An LCD panel beneath a transparent imaging surface \cite{ElMasry2019} allows the background colour to be set per variety. A dual-illumination architecture combining a top-mounted RGB ring light with a variable LCD background became the basis for subsequent prototypes.

\begin{figure}[H]
    \centering
    \begin{subfigure}[b]{0.30\textwidth}
        \centering
        \includegraphics[width=\textwidth]{BY - yellow from bottom only - for segmentation.jpg}
        \caption{BY}
        \label{fig:colour_by}
    \end{subfigure}
    \hfill
    \begin{subfigure}[b]{0.30\textwidth}
        \centering
        \includegraphics[width=\textwidth]{WB - white from top only - for classiication.jpg}
        \caption{WB}
        \label{fig:colour_wb}
    \end{subfigure}
    \hfill
    \begin{subfigure}[b]{0.30\textwidth}
        \centering
        \includegraphics[width=\textwidth]{BR - Red from bottom only - for multispectral classification.jpg}
        \caption{BR}
        \label{fig:colour_br}
    \end{subfigure}
    \caption{Final illumination configurations: (a)~BY yellow backlight for grain boundary segmentation, (b)~WB white top-light for general classification, and (c)~BR red backlight for multispectral defect classification.}
    \label{fig:figure4_1}
\end{figure}

%----------------------------------------------------------------------------------------
\section{Prototype I: Low-Fidelity Proof of Concept}
%----------------------------------------------------------------------------------------

Prototype I used a 3D-printed enclosure, a small LCD panel as the active background, WS2812B top illumination, and a Raspberry Pi 5 for control and capture. The objective was to verify the dual-illumination hypothesis in practice. The combined top-light and programmable backlight improved visual separation between categories that static backgrounds could not distinguish. Early-stage \textit{Infested} specimens became more distinguishable under certain background settings. However, the small enclosure handled only a fraction of a 100-gram sample. Scaling to operational volume required a separate design effort.

\begin{figure}[H]
    \centering
    \includegraphics[width=0.65\textwidth]{3d printed prototype.png}
    \caption{Prototype~I: low-fidelity 3D-printed enclosure validating the dual-illumination hypothesis.}
    \label{fig:figure4_2}
\end{figure}

%----------------------------------------------------------------------------------------
\section{Prototype II: Scaling to Industrial Volume and the Density Problem}
%----------------------------------------------------------------------------------------

Prototype II replaced the 3D-printed chassis with aluminium extrusion for durability and dust tolerance. The imaging surface was enlarged to accommodate a full 100-gram sample. Under full sample loading, a 100-gram wheat sample does not spread into a monolayer. At approximately 2,800 entities on a flat surface, grains stack against each other. Many specimens are entirely hidden from the overhead camera at any given moment. Inference on Prototype II images produced undercounting, merged detections, and poor classification confidence on occluded specimens. The constraint was in the imaging geometry, not the algorithm.

\begin{figure}[H]
    \centering
    \begin{subfigure}[b]{0.48\textwidth}
        \centering
        \includegraphics[width=\textwidth]{aluminium prototype.jpg}
        \caption{Aluminium extrusion chassis of Prototype~II.}
        \label{fig:aluminium_chassis}
    \end{subfigure}
    \hfill
    \begin{subfigure}[b]{0.48\textwidth}
        \centering
        \includegraphics[width=\textwidth]{aluminium prototype sample image.jpg}
        \caption{100-gram sample exhibiting high-density clumping.}
        \label{fig:aluminium_sample}
    \end{subfigure}
    \caption{Prototype~II at industrial scale, showing the high-density occlusion problem that redirected the design brief.}
    \label{fig:figure4_3}
\end{figure}

%----------------------------------------------------------------------------------------
\section{The Hardware-Software Co-Design Pivot}
%----------------------------------------------------------------------------------------

After Prototype II, the design question shifted from ``how do we detect occluded grains?'' to ``how do we ensure grains are not occluded during capture?''

This reframing follows the hardware-software co-design principle \cite{Keutzer2000}: rather than asking software to compensate for physical limitations, the hardware is redesigned so that the computational task becomes tractable. In this case, a mechanical system was required that could disperse a 100-gram sample across a sequence of frames while varying specimen orientations.

The dispersion mechanism had to meet four requirements: repeatability across hundreds of cycles, no damage to grain, compatibility with fine agricultural dust, and a total cycle time under five minutes. Chapter~\ref{Chapter5} describes the mechanical development that met these requirements.
